\chapter{Vitalfunktionen --- \gAasRsN{RS~V}}


Sebastian Gabriel [R]; Michael Withofner

\gAuthNotes{
Changelog:

\begin{description}
    \item[2009-02-25] Start Changelog. Kapitel mit Korr
    Koch korr. 
    \item[2009-mm-dd] in \LaTeX übernommen
    \item[2009-04-19] anhand der OpenOffice-Version überarbeitet.
Bilder fehlen noch. 
\item[2009-05-06] GAB: Abschnitte "`Blutdruck"' und
"`Atemvolumina"' aus Anatomie-Teil übernommen.
\end{description}

}

\minitoc

\paragraph{\"Uberblick}

\begin{center}
\begin{tabularx}{\linewidth}[H]{XX}
\toprule
Vitalfunktionen I. Ordnung
 &
Vitalfunktionen II. Ordnung
\\\midrule
\begin{itemize}
\item Bewu{\ss}tsein
\item Atmung
\item Kreislauf
\end{itemize}
\gFazit{Wichtigste, grundlegende Funktionen des K\"orpers

Ausfall einer dieser Funktionen bedeutet immer akute
Lebensgefahr und mu{\ss} dringlich behandelt werden!}
 &
\begin{itemize}
\item Wasser-/Elektrolythaushalt
\item S\"aure-/Basenhaushalt
\item Temperaturhaushalt
\item Hormonsystem
\item Stoffwechsel
\item Immunsystem
\end{itemize}
\\\bottomrule
\end{tabularx}
\end{center}


\paragraph[NACA -- Score]{NACA -- Score}
\gZ{\index{NACA}National Advisory Committee for Aeronautics

Diente urspr\"unglich der Bewertung von  Erkrankungen / Verletzungen und
der Entscheidung welche Soldaten zuerst auszufliegen w\"aren. Heute
dient der NACA-Score der groben statistischen Erfassung.  
}
1 bis 7:


\begin{table}[h]\begin{center}
\begin{tabular}{cl}\toprule
  NACA & Bedeutung
\\\midrule
 I & keine \"arztliche Therapie erforderlich
\\
 II & ambulante Abkl\"arung
\\
 III & station\"are Abkl\"arung
\\
 IV & akute Lebensgefahr nicht ausgeschlossen
\\
 V & akute Lebensgefahr
\\
 VI & Reanimation
\\
 VII & Tod
\\\bottomrule
\end{tabular}
\captionof{table}{NACA-Score}
\end{center}\end{table}


\section{Bewusstsein}
\gDefinition{\"Uberbegriff u.a. f\"ur Wachheit,
Orientierung, Aufmerksamkeit, Auffassungsgabe, Denkverlauf und
Merkf\"ahigkeit\label{ref:JRcite1Pschyrembel259}[1].} 

\gZ{Die {\quotedblbase}wahre{\textquotedblleft} Definition
von Bewusstsein ist seit tausenden von Jahren ungekl\"art. Da f\"ur die
Sanit\"aterausbildung aber nur beschr\"ankt Zeit zur Verf\"ugung steht
ist die obige Definition f\"ur den Alltag erstmal
ausreichend.}

Das Bewusstsein ist besonders wichtig hinsichtlich des Schutzes des
Menschen vor Bedrohung. Ein bewusstseinsklarer Mensch kann sich gegen
Gefahren wehren, ein eingetr\"ubter oder bewusstloser Mensch kann dies
schlecht oder gar nicht. 

\subsection{Bewu{\ss}seinsst\"orung}

\gDefinition{St\"orung des Bewusstseins, besonders der
Wachheit und Orientierung.}

\paragraph{Orientierung}

Es ist wichig zu \"Uberpr\"ufen ob der Patient orientiert ist. Man
unterscheidet dabei: 

\begin{enumerate}
    \item Orientierung zur Person: Name, Alter.
    \item \"Ortliche Orientierung: Wo befindet sich der Patient gerade
    (Stadt, Land, Adresse, Wohnung, Auto, etc.)?
    \item Zeitliche Orientierung: Tageszeit, Wochentag, Monat, Jahreszeit,
    Jahr?
\end{enumerate}

\paragraph{Allgemein}

Wie bereits oben erw\"ahnt ist das Bewusstsein wesentlich an der Abwehr
von Gefahren beteiligt. Je mehr das Bewusstsein gest\"ort oder
reduziert ist desto gef\"ahrlicher ist das f\"ur den Patienten. 

\begin{itemize}
    \item Ab einem gewissen Punkt sind auch so wesentliche Funktionen wie
    der \gT{Schluck-}, \gT{W\"urge-} bzw. \gT{Hustenreflex}
    ausfallen und es kann zum verschlucken von Mageninhalt oder Blut in die Lunge kommen
    (\gT{Aspiration}). Man spricht auch vom \gEs{Ausfall der
    Schutzreflexe}.
    \item Bei schweren Bewusstseinsst\"orungen kann es auch zum Erschlaffen
    wichtiger Muskeln wie z.B. Der Zunge kommen. F\"allt diese zur\"uck
    kann sie den Atemweg blockieren oder ganz verschlie{\ss}en.
\end{itemize}

\gFazit{Daher sind schwere Bewusstseinsst\"orungen
\gEs{lebensgef\"ahrlich!}}

\paragraph{Ursachen}

${\rightarrow}$ Sch\"adigung des zentralen Nervensystems

\begin{center}
\begin{tabularx}{\linewidth}[H]{XX}
\toprule
\gEs{Prim\"are Ursachen} & \gEs{Sekund\"are Ursachen}
\\
{\tiny betreffen direkt das ZNS} & {\tiny betreffen \"uber
Umwege das ZNS} 
\\\midrule
\begin{itemize}
    \item Verletzungen (SHT)
    \item Blutungen, Hirnschwellung
    \item Schlaganfall
    \item epileptischer Anfall
    \item Entz\"undungen
    \item \gZ{Tumor}
\end{itemize}
$\vdots$
 &
 \begin{itemize}
     \item St\"orungen der $O_2$-Versorgung:
     \begin{itemize}
         \item Atemst\"orungen  (Hypoxie)
         \item Herz-Kreislaufst\"orungen
     \end{itemize}
     \item Störungen des Wasserhaushaltes und
     Elektrolytentgleisungen
     \item Stoffwechselst\"orungen
     \begin{itemize}
         \item Zuckerhaushalt
         \item Leberversagen
     \end{itemize}
     \item Vergiftungen:
     \begin{itemize}
         \item Alkohol
         \item Medikamente, Drogen
         \item Reizgase,  L\"osungsmittel ...
     \end{itemize}
 \end{itemize}
$\vdots$
\\\bottomrule
\end{tabularx}
\end{center}

\subsubsection{Grade und Symptome}

Symptome Bewusstseinsst\"orungen: Tabelle geh\"ort erg\"anzt.%
%



\begin{table}[H]\begin{center}
\begin{tabulary}{\linewidth}{LLLL}
\toprule
Bewu{\ss}tseinsklar
 &
 &
 &
\\\midrule
Bewu{\ss}tseinsgetr\"ubt
&
somnolent
 &
Schl\"afrig, aber leicht erweckbar.
 &
\\
 &
sopor\"os
 &
 &
Gefahr!
\\\midrule
Bewu{\ss}tlos
 &
Koma
 &
 &
Lebensgefahr!
\\\bottomrule
\end{tabulary}
\caption{Wachheitszustände}
\end{center}\end{table}


\paragraph{Gefahren}

\begin{itemize}
    \item Zur\"uckfallen der Zunge ${\rightarrow}$ Erstickung!
    
    \item \gEs{Aspiration}  (Husten- bzw. W\"urgereflex
    sowie s\"amtliche andere Schutzreflexe ausgefallen,  Muskulatur komplett erschlafft
    ${\rightarrow}$  Erbrochenes gelangt ungehindert in die  Lunge!)
    
\end{itemize}

\paragraph{Glasgow Coma Scale (GCS)}\index{Glasgow Coma
Scale}\index{GCS}

Die Glasgow Coma Scale dient der objektiven Einsch\"atzung
der Bewu{\ss}tseinslage mittels einer Skala mit 3 bis max.
15 Punkten.

Bewu{\ss}tseinsklarer Patient: 15 Punkte. Toter Patient: 3 Punkte (nicht
0!)

Bewertet werden: 


\begin{table}[H]\begin{center}
\begin{tabulary}{\linewidth}{LCCL}
\toprule
~ & Max. & Punkte & Reaktion
\\\midrule
Augen\"offnung & 4 & 4 & spontan
\\
 & & 3 & auf Ansprache
\\
 & & 2 & Schmerzreiz
\\
 & & 1 & nicht
\\\midrule
Verbale Antwort & 5 & 5 & klar
\\
 & & 4 & verwaschen
\\
 & & 3 & stammelnd
\\
 & & 2 & unartikuliert
\\
 & & 1 & nicht
\\\midrule
Motorische Reaktion & 6 & 6 & spontan
\\
 & & 5 & gezielte Abwehrbewegungen
\\
 & & 4 & ungezielte Abwehrbewegungen
\\
 & & 3 & Beugung
\\
 & & 2 & Strecken
\\
 & & 1  & bewegt gar nicht
\\\bottomrule
\end{tabulary}\end{center}
\caption{GCS}
\end{table}



\begin{gMassnahmenEnv}
    \paragraph{Ma{\ss}nahmen bei getr\"ubten Patienten}

    Ma{\ss}nahmen bei Bewusstseinsst\"orungen: Neufassung, bitte besonders
    genau durchsehen!%

    \begin{itemize}
        \item BAK -- Schema  (Bewu{\ss}tsein/Atmung/Kreislauf)
        \item Lagerung: situationsgerecht je nach
        Verdachtsdiagnose, im Zweifel \gEs{stabile
        Seitenlage}, bei hochschwangeren Frauen auf die
        \gE{linke} Seite.
        \item \gTxBeiVit \gTxMassVitBes

        \begin{itemize}
            \item $O_2$: nach Verdachtsdiagnose, im Zweifel
            ca. 6-8\gLMin
        \end{itemize}

        \item Engmaschiges \gEs{Monitoring} (Patient
        beobachten und regelmäßig ansprechen, Atmung,
        Pulsoxymetrie, RR)
        \item W\"armeerhalt
        \item \gEs{Ursachenforschung} -- Warum ist der
        Patient eingetr\"ubt?
        \item \gEs{Neurcheck}
        \begin{itemize}
            \item Immer auch an \gEs{BZ} denken!
        \end{itemize}
        \item \gEs{Traumacheck}. Suche nach Primär- und
        Folgeverletzungen!
    \end{itemize}

\end{gMassnahmenEnv}

\begin{gMassnahmenEnv}
    \paragraph{Ma{\ss}nahmen bei bewu{\ss}tlosen und sopor\"osen Patienten}
    \begin{itemize}
        \item BAK -- Schema  (Bewu{\ss}tsein/Atmung/Kreislauf)
        \item \gTxMassVitBes
        \begin{itemize}
            \item Lagerung: stabile Seitenlage, bei hochschwangeren Frauen auf die
        \gE{linke} Seite. 
            \item $O_2$:  10\gLMin , bei
            Zyanose\footnote{Zyanose: Bl\"auliche Verf\"arbung der Haut infolge einer unzureichenden
            Sauerstoffversorgung. Besonders auff\"allig an den Lippen, Gesicht, und
            Extremit\"aten.}  15\gLMin
        \end{itemize}
        \item W\"armeerhalt
        \item \gEs{Ursachenforschung} -- Warum ist der Patient bewusstlos oder
        sopor\"os?
        \item \gEs{Neurcheck}
        \begin{itemize}
            \item Immer auch an \gEs{BZ} denken!
        \end{itemize}
        \item \gEs{Traumacheck}. Suche nach Primär- und
        Folgeverletzungen!
    \end{itemize}
\end{gMassnahmenEnv}



\section{Atmung}

\minisec{Kurze Wiederholung}

normale Atemfrequenz (\gT{AF}) (Erwachsener): 12-15x / min,
Atemzugvolumen (\gT{AZV}) $\approx$ 5-7 ml\,/\,kg\,KG

Atemminutenvolumen (AMV): $AMV = AZV \times AF$

\index{Totraumventilation}\gT{Totraumventilation}: Luftmenge
in oberen  und teilweise unteren Atemwegen, die  nicht am Gasaustausch teilnimmt.
Achtung \index{Schnappatmung}\gT{Schnappatmung}! Dabei wird
nur so wenig Luft bewegt sodass aufgrund des Totraumes keine frische Luft in die
Lungenbl\"aschen (Alveolen) kommt. H\"aufig wird die Schnappatmung mit
flacher Atmung verwechselt, aber:

\gCave{Schnappatmung = keine Atmung}


\subsubsection{Atemvolumina}\label{topic:atemvolumina}
\gAuthNotes{KOCH: zu Vitalfunktionen}


\begin{center}

\gTempPicMissing{Atemvolumina img43}

\end{center}
+ 150ml in den oberen Luftwegen (Trachea, Rachenraum) ${\rightarrow}$
Dieser Teil nimmt nicht am Gasaustausch ${\rightarrow}$
\index{Totraum}\gT{Totraum} 

normale Atemfrequenzen sind altersabh\"angig:

\begin{center}
\begin{tabular}{lrr}
\toprule
 & AF &  AZV
\\\midrule
Erwachsener & 12 -- 16/min  & ca. 500 ml
\\
Neugeborenes & 40/min & 20-40 ml 
\\
Kleinkind & 25/min & 100-200 ml
\\
Schulkind & 20/min & 200-400 ml
\\\bottomrule
\end{tabular}
\end{center}

\gFazit{AMV = AZV x AF, beim Erwachsenen  ca. 6-8\l}

AF  ... Atemfrequenz

AMV  ... Atemminutenvolumen

AZV  ... Atemzugsvolumen ... Volumen bei einem normalen Atemzug

\section{St\"orungen der Atmung}

\gT{Ateminsuffizienz} \ldots eingeschr\"ankte Atmung
$\rightarrow$ $O_2$-Mangel (Hypoxie)   

\gT{Atemstillstand} \ldots Aussetzen d. Atmung \gZ{(Apnoe)}

\paragraph{Ursachen}~

\begin{tabulary}{\linewidth}[H]{LL}
\toprule
St\"orung des/der  &
\\\midrule
$O_2$-Angebotes
 &
G\"arkeller ($CO_2$\,${\uparrow}$, $O_2$\,${\downarrow}$)

{\quotedblbase}d\"unne Luft{\textquotedblleft} (z.B. Hochgebirge,

p$O_2$\,${\downarrow}$)

Defekte Gastherme ($CO$\,${\uparrow}$ $O_2$\,${\downarrow}$)

Ertrinken
\\\midrule
 $O_2$-Diffusion
 &
Lungen\"odem

Lungenentz\"undung

Lungenembolie
\\\midrule
 Atemmechanik
 &
Verlegung der oberen Atemwege (Zunge, Erbrochenes, Laryngospasmus,
         Glottis\"odem, Bolus)

Verlegung der unteren Atemwege (Bronchitis, Asthma,
Lungen\"odem)

Verminderung der Dehnbarkeit der Thoraxwand oder des Lungengewebes
     (Rippenfraktur, Pneumothorax, Pleuraergu{\ss}, Emphysem)
\\\midrule
 Neuromuskul\"aren Atemregulation
 &
SHT (Sch\"adel Hirn Trauma)

Schlaganfall (Insult)

Vergiftungen

Tumore

hoher Querschnitt (RM durchtrennt)

Muskelerkrankungen

Nervenerkankung

Stoffwechselst\"orungen
\\\bottomrule
\end{tabulary}


\paragraph{Symptome}~

\begin{center}
\begin{table}[H]
\begin{tabulary}{\linewidth}[H]{LLLL}
\toprule
Gruppe
 &
Befund
 &
Beschreibung
 &
Anmerkung
\\\midrule
Dyspnoe
 &
Atemnot
 &
Leitsymptom
 &
\\\midrule
Atemger\"ausch
 &
Brodeln
 &
Blubbern, klassisch f\"ur Lunden\"odem
 &
\\
 &
Stridor
 &
 &
\\
 &
Brummen, Giemen
 &
 &
\\
 &
Rasselger\"ausche
 &
 &
\\
Frequenz
 &
Beschleunigt
 &
Tachypnoe
 &
\\
 &
Verlangsamt
 &
Bradypnoe
 &
\\\midrule
Atemzugsvolumen
 &
Schnappatmung
 &
 &
\\
 &
Flache Atmung
 &
 &
\\
 &
Tiefe Atmung
 &
 &
\\\midrule
Hautfarbe
 &
Blass
 &
Normal, An\"amie, Blutverlust
 &
\\
 &
Rosig
 &
Normal, CO-Vergiftung
 &
\\
 &
Bl\"aulich 
 &
Zyanose: Hypoxie!
 &
\\\midrule
K\"orperlich
 &
Einziehungen an den Rippen
 &
Einsatz der Atemhilfsmuskulatur
 &
\\
 &
Aufst\"utzen
 &
Einsatz der Atemhilfsmuskulatur
 &
\\
 &
Nasenfl\"ugeln
 &
Atemnot, besonders bei Kleinkindern
 &
\\
 &
Aufrechte Position
 &
 &
\\
 &
{\quotedblbase}Paradoxe Atmung{\textquotedblleft}
 &
Brustkorb senkt sich bei Einatmung
 &
\\\bottomrule
\end{tabulary}
\caption{Symptome: Atmung}
\end{table}
\end{center}


\paragraph{Gefahr der Ateminsuffizienz}

\begin{itemize}
\item Zu geringer $O_2$-Gehalt im Blut (Hypoxie)
\begin{itemize}
\item in weiterer Folge Unterversorgung lebenswichtiger Organe wie Herz,
Hirn.
\end{itemize}
\item Atmungsbedingte \"Ubers\"auerung (Respiratorische
Azidose, s. \ref{topic:azidose-respiratorische} )
\end{itemize}

\begin{gMassnahmenEnv}
    \paragraph{Spezielle Ma{\ss}nahmen bei Atemstillstand}

\begin{itemize}
    \item Freimachen der Atemwege
    \begin{itemize}
        \item \"Offnen des Mundes (Esmarch-Handgriff)
        \item Inspektion
        \item fals n\"otig ausr\"aumen
        \item (falls n\"otig) Absaugen
        \item \"Uberstrecken der HWS
        \item Zungengrund wird angehoben -{\textgreater} Atemwege  sind frei!
    \end{itemize}
    \item \"Uberpr\"ufung der Atemt\"atigkeit (bei  \"uberstreckter HWS)
    durch
    \begin{itemize}
        \item SEHEN, H\"OREN, F\"UHLEN f\"ur mindestens 10  Sekunden!
    \end{itemize}
    \item Sauerstoffgabe (wenn ausreichende Spontan-Atmung vorhanden)
    \item Beatmung (keine ausreichende Spontanatmung bzw. bei Reanimation)
    \begin{itemize}
        \item Beatmungsbeutel mit Maske und Reservoir
    \end{itemize}
\end{itemize}

Vgl. die entsprechenden Kapitel im Teil Erste Hilfe und
Sanit\"atstechnik bez\"uglich der genauen Durchf\"uhrung der
Ma{\ss}nahmen. 
\gCave{Achtung: Bei einem mit Atemstillstand vorgefundenen
Patienten ist grunds\"atzlich von einer Reanimation
auszugehen.}
\end{gMassnahmenEnv}









\section[Herz- / Kreislauf]{Herz- / Kreislauf}
\label{bkm:RefHeading42514202}St\"orungen des Herz-/ Kreislaufsystems



Durch die Pumpfunktion des Herzens ensteht ein
lebenswichtiger Druck im (arteriellen) Gefäßsystem, der
\gE{Blutdruck}. Vgl. Kap. \emph{Vitalfunktionen},
\ref{topic:blutdruck}. Abk\"urzung: RR. Gemessen immer in Herzh\"ohe!

\subsection{Blutdruck} 

Durch die Pumpfunktion des Herzens ensteht ein
lebenswichtiger Druck im (arteriellen) Gefäßsystem, der
\gE{Blutdruck}.

\gMarried
{Durch die rhythmische Funktion des
Herzens kommt es infolge des Blutauswurfes zu genauso rhythmischen Druckveränderungen.
Während des Auswurfes von Blut in das
Gefäßsystem (Herzanspannung, Systole)
spricht man vom systolischen Blutdruck, während
der Enstpannungsphase des Herzens (Diastole) vom
diastolischen Blutdruck. Nattürlich ist der Blutdruck bei
der Anspannung des Herzens höher als bei der Entspannung
(Systolischer Blutdruck > Diastolische Blutdruck).}
{\begin{description}
    \item[systolisch:] Maximaldruck in der Arterie am H\"ohepunkt der
Kammerkontraktion
    \item[diastolisch:] Druck in der Arterie w\"ahrend die Kammer erschlafft ist
\end{description}}







Normbereich : zwischen 90/60 mmHg und 140/90 mmHg  

Hypotonie:  {\textless} 90/60 mmHg

Hypertonie: {\textgreater} 140/90mmHg  

Achtung: Eine Hypertonie ist zwar eine Erkrankung, aber nicht immer im
Einsatz von Bedeutung. Viele Patienten sind an ihren erh\"ohten
Blutdruck gewohnt und haben keine Beschwerden. Es darf nie nur nach dem
Wert gehandelt werden, es muss auch immer bedacht werden: 

\begin{itemize}
\item Zustand des Patienten.

\item Der sonst \"ubliche Blutdruck des Patienten (erfragen,
Blutdruckprotokoll).

\item Der Verlauf des Blutdruckes w\"ahrend der Behandlung
(Traumapatienten !!!).

\item Der Erregungszustand des Patienten.

\end{itemize}

\minisec{Beispiele}

\emph{Ein 65-j\"ahriger, m\"annlicher Patient klagt \"uber
Schwindelgef\"uhl, sie messen einen Blutdruck von 100/70. Im Blutdruckprotokoll finden sie
f\"ur die vergangenen vier Tage folgende Eintr\"age: 170/110, 180/115,
175/113, 183/112, ... .}


Wie beurteilen Sie den Blutdruck? 

\begin{quote}
Der Blutdruck ist f\"ur den Patienten zu niedrig, der K\"orper ist auf
eine so pl\"otzliche Umstellung nicht vorbereitet gewesen. Im
Anamnesegespr\"ach erfahren sie dass er heute das erste Mal einen
{\quotedblbase}Nitro-Spray{\textquotedblleft} verwendet hat. Sie
vermerken dies am Einsatzprotokoll. Im Krankenhaus erfahren Sie dass
wahrscheinlich eine \"Uberdosierung dieses Medikaments die Ursache
dieses Blutdruckabfalls war. ${\rightarrow}$ Obwohl der Patient laut
Lehrbuch einen {\quotedblbase}sch\"onen{\textquotedblleft} Blutdruck
hatte, war dieser dennoch zu niedrig!
\end{quote}

\emph{Ein 56-j\"ahriger Patient klagt \"uber starke
kolikartige Schmerzen im linken Unterbauch. Sie messen einen Blutdruck von 155/95. Der Patient
kennt seinen {\quotedblbase}gew\"ohnlichen{\textquotedblleft} Blutdruck
nicht.}  

Wie beurteilen Sie den Blutdruck? 

\begin{quote}
Der Blutdruck ist laut Lehrbuch erh\"oht, aber dies ist erkl\"arbar die
starken Schmerzen und den damit verbundenen Erregungszustand. Die
Beschwerden stehen wahrscheinlich in keinem urs\"achlichen Zusammenhang
mit den Schmerzen. 
\end{quote}

\emph{Eine 50-j\"ahrige Patientin hat Nasenbluten. Sie
messen einen Blutdruck von 220/120. Auf Nachfrage gibt sie an etwas Kopfsschmerzen zu haben,
und etwas schwindlig sei sie auch.} 

Wie beurteilen Sie den Blutdruck? 

\begin{quote}
Der Blutdruck ist sehr stark erh\"oht, ausserdem zeigt die Patientin
Zeichen einer hypertensiven Krise (Nasenbluten, Kopfweh, Schwindel).
Hier ist der erh\"ohte Blutdruck Ursache f\"ur die Beschwerden und muss
mittels Medikamente unbedingt gesenkt werden. 
\end{quote}



\subsection{Kreislaufstillstand}\label{bkm:RefHeading42514602}

\gDefinition{Totales Pumpversagen des Herzens} aufgrund von

\begin{itemize}
    \item Asystolie (keine Herzt\"atigkeit)
    \item Kammerflimmern (unkoordinierte Herzt\"atigkeit)
    \item Zu schnelle Herzfrequenz, bei der sich das Herz
    nicht mehr füllen oder etwas auswerfen kann ((Pulslose) Ventrikul\"are Tachykardie)
    \item Elektromechanische Entkopplung: keine
    Pumpreaktion auf Signale des Reizleitungssystems im Herz
\end{itemize}


\paragraph{Ursachen}

\begin{itemize}
    \item Minderversorgung des Herzmuskels mit O2
    \begin{itemize}
        \item Herzinfarkt
        \item Herzinsuffizienz
    \end{itemize}
    \item St\"orungen der Atmung
    \begin{itemize}
        \item Asthma bronchiale
        \item Lungenembolie
        \item zentrale Atemregulationsst\"orungen
    \end{itemize}
    \item sonstige St\"orungen
    \begin{itemize}
        \item Trauma
        \item Vergiftungen
        \item thermische Sch\"aden
        \item Elektrolytverschiebungen,...
    \end{itemize}
\end{itemize}


\paragraph{Pathophysiologie}

\begin{itemize}
\item Aussetzen des Herzschlages $\rightarrow$ Patient ist
sofort pulslos!
\item Bewu{\ss}tlosigkeit (nach 15-20\,sek)
\item Atemstillstand/Schnappatmung (nach ca. 10\,sek)
\item \gEs{Irreversible Hirnsch\"aden nach 3-5\,min}, bei
K\"alte allerdings deutlich l\"anger!
\end{itemize}











\subsection{Herzinsuffizienz, Herzversagen}

Weitere, oft auch sehr bedrohliche, St\"orungen des Herzens werden im
Kapitel  Herz-Kreislaufst\"orungen 
(\ref{bkm:RefHeading42342802}
auf S.\,\pageref{bkm:RefHeading42342802}:
{\quotedblbase}Herz-Kreislaufst\"orungenunten{\quotedblbase}) besprochen.

\section{SCHOCK}

\gDefinition{{\quotedblbase}Lebensbedrohliche Kreislaufst\"orung aufgrund von
Unterversorgung des K\"orpers mit Blut und Sauerstoff infolge des
Versagens des Kreislaufsystems.{\textquotedblleft}}

\gEs{Es k\"onnen einzelne oder mehrere Teile des Kreislaufsystems betroffen
sein  (Blut, Gef\"a{\ss}e, Herz).}

\gEs{Ob ein Schock oder eine Schockgefahr besteht wird anhand von
\underline{Schocksymptomen} und anhand der
\underline{Verdachtsdiagnose} beurteilt.}


Schock bedeutet:

\begin{itemize}
    \item akute Kreislaufinsuffizienz = 
   \gEs{Mi{\ss}verh\"altnis
   Herzauswurf\,:\,Sauerstoffbedarf} des K\"orpers
    \item M\"ogliche Ursachen:
    \begin{enumerate}
        \item Herz pumpt wenig
        \item Zu wenig Blut vorhanden
        \item Gef\"a{\ss}e sind zu weit gestellt und Blut versackt
    \end{enumerate}
\end{itemize}

\gFazit{Im Endeffekt kommt immer zu wenig Blut mit
Sauerstoff zu den Organen, dadurch ist die Ge\-webedurchblutung nicht ausreichend  f\"ur die
Versorgung aller Organe mit Sauer\-stoff.} 

K\"orpereigene Mechanismen k\"onnen gegensteuern und versuchen den
Schock zu kompensieren (gegensteuern). Dies geht allerdings nur eine
gewisse Zeit und in einem gewissen Ausma{\ss} gut, daher ist eine
schnell  einsetzende Therapie wichtig! Bei l\"anger bestehendem Schock
kann dieser Zustand irreversibel werden. 

\gE{Kein} Schock besteht bei kurzzeitigen 
Bewu{\ss}tseinsst\"orungen wie Kollaps/Synkope
(\underline{kurzfristige} Blutverteilungsst\"orung,
kurzfristige Minderdurchblutung des Gehirns,  fl\"uchtige Symptome)


\paragraph{Schock --  Kompensationsmechanismen}

Der K\"orper versucht selber dagegen etwas zu tun und gegenzusteuern
({\quotedblbase}kompensieren{\textquotedblleft}). Der wichtigste
Mechanismus ist die \index{Zentralisation}\gT{Zentralisation}. Darunter versteht man die Blutumverteilung zugunsten lebenswichtger Organe wie
Herz, Gehirn, Lunge und Leber. Dabei werden z.B. Haut, Muskel und Darm
sowie die Extremit\"aten (=Peripherie) kaum mehr  durchblutet.
Periphere Venen kollabieren und die Rekapillarisierungszeit
verl\"angert. Wenn die Kompensationsma{\ss}nahmen  versagen spricht man
vom \gE{"`dekompensierter
Schock"'}. Es besteht sp\"atestens dann akute
\gEs{Lebensgefahr}!

Beachte: 

\gCave{Bei jedem Patienten ist zu erheben ob ein Schock oder eine realistische
Schockgefahr besteht oder nicht bzw. festzustellen um welche Art von
Schock es sich handelt (s.u.)!}

\gCave{Achtung! 

Jeder Schockzustand ist grunds\"atzlich ein
Fall f\"ur den NOTARZT!}

\subsection{Allgemeine Schocksymptomatik}

\begin{center}
\begin{tabular}{|m{3.9750001cm}m{9.325cm}|}
\hline
 &
Symptome
\\\hline
{\sffamily\Large\color{red} \index{Schockgefahr}Schockgefahr}

\centering\includegraphics[width=1.5cm]{./icons/sm-basic.png}
 &
\begin{itemize}
\item Keine eindeutigen Anzeichen

\item Vorausschauend handeln!

\item Bedenke Verdachtsdiagnose.  Muss ich damit rechnen dass sich ein
Schock entwickelt?  Besteht Schockgefahr?

\end{itemize}
\\\hline
{\sffamily\Large\color{red}Schock}

\centering\includegraphics[width=1.5cm]{./icons/sm-pale.png}
 &
\begin{itemize}
\item HF ${\uparrow}$

\item m\"a{\ss}iger RR-Abfall

\item RR noch {\textgreater} HF

\item neurologisch unauff\"allig, ev. Euphorisch oder agitiert

\item Haut bla{\ss}, evt. kaltschwei{\ss}ig

\item Venen sind kollabiert oder gestaut

\item Peripherie schlecht durchblutet

\end{itemize}
\\\hline
{\sffamily\Large\color{red}Schwerer Schock

Vollbild}

\centering\includegraphics[width=1.5cm]{./icons/sm-pale-frown.png}

 &
\begin{itemize}
\item Tachykardie (HF ${\uparrow}$,
{\quotedblbase}fliehender{\textquotedblleft} Puls)

\item RR-Abfall (schwacher Puls)

\end{itemize}
\begin{itemize}
\item Herzfrequenz {\textgreater} systolischer RR

\end{itemize}
\begin{itemize}
\item Tachypnoe (schnelle, flache Atmung)

\item neurologisch auff\"allig, getr\"ubt

\item Haut bla{\ss}, zyanotisch (bl\"aulich), feucht

\item Venen kollabiert oder gestaut

\item nur mehr schwer behebbar! ${\rightarrow}$ Schock  m\"oglichst
fr\"uhzeitig erkennen \&  behandeln

\end{itemize}
\\\hline
{\sffamily\Large\color{red}Endstadium}

\huge\centering $\skull$
 &
\begin{itemize}
\item Haut grau, klebrig

\item Versagen der Herzfunktion

\item RR kaum mehr me{\ss}bar

\item evt. Koma

\item evt. Schnappatmung

\item Patient ist dabei zu sterben.

\end{itemize}
\\\hline
\end{tabular}
\end{center}

\paragraph{Schockschere}

Normalerweise ist die Herzfrequenz niedriger als der systolische
Blutdruck. Wenn der K\"orper versucht den Schock zu kompensieren kann
sich dieses Verh\"altnis umkehren. Man spricht dann von der
Schockschere, der Patient befindet sich dann bereits schon im Vollbild
eines Schocks! 

Achtung: \gEs{Kinder haben andere Normalwerte f\"ur HF und
RR!} Die Faustregel ist deshalb nicht anwendbar!

\captionof{figure}{Schockschere. \gAuthNotes{Achtung! Bild
fehlt!}}

\gFazit{Wenn sich das Verh\"altnis Herzfrequenz zu Systolischem
Blutdruck umkehrt befindet sich der Patient bereits im Schock.}

\subsection{Allgemeine Ma{\ss}nahmen der
Schockbek\"ampfung}\marginlabel{Schockbekämpfung}\label{topic:schockbekaempfung}\index{Schockbek\"ampfung}
\label{bkm:RefHeading38104902}

\gFazit{Grundsatz: \gEs{Der Schock muss bek\"ampft werden
bevor man ihn bemerkt.}}

\begin{gMassnahmenEnv}
    \begin{itemize}
        \item Notarzt!
        \item Situationsgerechte  Lagerung, je nach Diagnose und Schockart
        \item Beengende Kleidung \"offnen
        \item Sauerstoffgabe hochdosiert (ab 6\gLMin bei Schockgefahr,
        10\gLMin oder mehr bei bestehendem Schock, je nach Schwere)
        \item Monitoring (Pulsoxy, RR-Manschette angelegt lassen)
        \item Reanimationsbereitschaft
        \item W\"armeerhalt
        \item Psychische Betreuung
        \item Weitere Massnahmen je nach Schockart
    \end{itemize}
\end{gMassnahmenEnv}
    \gAuthNotes{Freigabe: Sauerstoffgabe hochdosiert (ab 6l/min bei
    Schockgefahr, 10l/min oder mehr bei bestehendem Schock, je nach Schwere}

\subsection{Schock -- Arten}

\begin{center}
\begin{tabularx}{\linewidth}[H]{XX}
\toprule
\gEs{Mechanismus}
 &
\gEs{Schockart}
\\\midrule
Herz pumpt zu wenig
&
{\begin{itemize}
    \item Kardiogender Schock
\end{itemize}}
\\
Zu wenig Blut vorhanden 
 &
{\begin{itemize}
    \item Hypovol\"amischer Schock
\end{itemize}}
\\
Gef\"a{\ss}e inad\"aquat weitgestellt, Blut versackt
 &
{\begin{itemize}
    \item Anaphylaktischer Schock
    \item Septischer Schock, toxischer Schock
    \item Neurogener Schock
\end{itemize}}
\\\bottomrule
\end{tabularx}
\end{center}


\subsubsection{Hypovol\"amischer
Schock}\label{topic:schock-hypovolaemischer}

\gDefinition{Schock aufgrund eines Verlustes von Blutvolumen aus den
Gef\"a{\ss}en.}

Es kann dabei verloren gehen: 

\begin{itemize}
    \item Blut
    \item Plasma
    \item Wasser und Elektrolyte
\end{itemize}

\paragraph{Ursachen}

\begin{center}
\begin{tabularx}{\linewidth}[H]{XX}
\toprule
Fl\"ussigkeitsverlust nach au{\ss}en
 &
Fl\"ussigkeitsverlust nach innen
\\\toprule
\begin{itemize}
    \item Blutung
    \item gastrointestinal (Durchfall, Diarrhoe )
    \item Verbrennungen (Hautbarriere versagt  gro{\ss}\-fl\"achig)
    \item renal (Diabetes mellitus, Diuretika, Polyurie bei  Nierenversagen)
\end{itemize}
 &
 \begin{itemize}
     \item Weichteilblutung
     \begin{itemize}
         \item Nach Trauma, rupturierte Eileiter\-schwanger\-schaft, etc
     \end{itemize}
     \item Wassereinlagerung ins Gewebe
     \item \gZ{Aszites
     ({\quotedblbase}Wasserbauch{\textquotedblleft} bei
     Lebererkrankungen  etc.)}
     \item Ileus (Darmverschlu{\ss})
     \item H\"amatothorax, H\"amatoperiton\"aum
\end{itemize}
\\\bottomrule
\end{tabularx}
\end{center}
Bei akutem Fl\"ussigkeitsverlust kann ein Verlust von bis zu 20\% des
Blutvolumens  gut kompensiert werden. Auf gr\"o{\ss}ere Verluste  folgt
ein Absinken von Herzauswurf und RR!

\begin{gMassnahmenEnv}
    \paragraph{Spezielle Ma{\ss}nahmen beim Hypovolämischen
    Schock}
    
    \begin{itemize}
        \item Allgemein: \gEs{Allgemeine Ma{\ss}nahmen der
        Schockbek\"ampfung} (s.
        \ref{topic:schockbekaempfung})
        \item Blutstillung  (weitere Verluste verhindern!)
        \item ggfs. Trauma-Vorgehen beachten
        \item Wenn m\"oglich schonende Bergung
        \item Lagerung: grunds\"atzlich Beine hoch
        
        Wichtige Ausnahmen: {\textperiodcentered}\,kardiale Ursache,
        {\textperiodcentered}\,Kopf- / Brust\-ver\-letzungen,
        {\textperiodcentered}\,Hirndruck (SHT),
        {\textperiodcentered}\,Knochenbr\"uche der Beine oder des
        Beckens, {\textperiodcentered}\,Bewusstlosigkeit
        ${\rightarrow}$ sonst Flachlagerung oder stabile Seitenlage
        \item evt. Aviso (Ank\"undigung) im  Krankenhaus / Schockraum
        \item Abtransport: rasch, schonend, nicht \"uberst\"urzt
    \end{itemize}
\end{gMassnahmenEnv}

\gAuthNotes{Begriff {\quotedblbase}Schocklagerung{\textquotedblleft}
ein f\"ur alle mal in der Ausbildung streichen, nur Hinweis dass dieser Ausdruck
veraltet ist sollten sie ihn wo aufschnappen. Stattdessen
{\quotedblbase}Beine-hoch{\textquotedblleft}-Lagerung.

Wenn KI gegen Beine hoch: Flachlagerung bzw. stab. Seitenlage empfehlen?
}


\subsubsection{Kardiogener Schock}
\gDefinition{Schock infolge mangelnder Herzpumpleistung.}
(Die Pulmonalembolie als Auslöser eines Schocks wird auch
zum kardiogenen Schock gerechnet, da es durch die
Verstopfung des kleinen Kreislaufes zu einer mangelhaften
Füllung des Herzens und damit zu einer mangelhaften
Auswurfleistung kommt.)

\begin{center}
\begin{tabularx}{\linewidth}[H]{XX}
\toprule
Ursachen &
charakterisiert durch\\\midrule
\begin{itemize}
    \item Herzversagen, Herzinsuffizienz
    \begin{itemize}
        \item Myokardinfarkt (h\"aufigste Ursache)
        \item Herzrhythmusst\"orungen
    \end{itemize}
    \item Lungenembolie
\end{itemize}
 &
 \begin{itemize}
     \item RR-Abfall
     \item eingeschr\"ankte Pumpleistung
     \begin{itemize}
         \item Blut staut sich zur\"uck
     \end{itemize}
 \end{itemize}
\\\bottomrule
\end{tabularx}
\end{center}

\begin{gMassnahmenEnv}
\paragraph{Spezielle Ma{\ss}nahmen beim Kardiogenen Schock}

\begin{itemize}
    \item Allgemein: \gEs{Allgemeine Ma{\ss}nahmen der
        Schockbek\"ampfung}
        (s.\,\ref{topic:schockbekaempfung})
    \item \gTxMassVitBes
    \begin{itemize}
        \item $O_2$: 10\gLMin   oder mehr
        \item Lagerung: Oberk\"orper hoch! Evt. Beine h\"angen  lassen.
    \end{itemize}
    \item Medikament\"ose Therapie durch Notarzt 
    \gZ{(Diuretika, Katecholamine, Heparin)}
\end{itemize}
\end{gMassnahmenEnv}







\subsubsection{Anaphylaktischer Schock}
\gDefinition{Schock aufgrund einer \"uberschie{\ss}enden
allergischen Reaktion -- {\quotedblbase}Allergischer
Schock{\textquotedblleft}.}

 Ausgel\"ost durch \"uberschie{\ss}ende Reaktion der K\"orerabwehr (Immunreaktion) auf
bestimmte Stoffe (Antigene). 

\begin{center}
 [Warning: Image ignored] % Unhandled or unsupported graphics:
%\includegraphics[width=2.641cm,height=3.837cm]{AASdev20090228-img70.jpg}
\gTempPicMissing{Allergischer Schock}
\end{center}
Es werden s\"amtliche Gef\"a{\ss}e maximal erweitert, zus\"atzlich
steigt die Durchl\"assigkeit der Gef\"a{\ss}e
(Gef\"a{\ss}permeabilit\"at ${\uparrow}$) ${\rightarrow}$ 
Plasmaverlust in den Zwischengebwebsraum (Interstitium)
${\rightarrow}$ RR f\"allt massiv ab!

\paragraph{M\"ogliche Antigene}

\begin{itemize}
\item Pollen
\item Medikamente, Antibiotika
\item Kontrastmittel
\item tierische Gifte (Bienenstich etc.)
\item und vieles andere ...
\end{itemize}

\paragraph{Symptome}

\begin{center}\captionof{table}{Stadien der Anaphylaxie}
\begin{tabulary}{\linewidth}[H]{cL}
\toprule
\gTabHeading{\mbox{Stadium}}
 &
\\\toprule
\gTabSub{I}
 &
Hautreaktion: evt. Ausschlag, R\"otung;
{\quotedblbase}Quaddeln{\textquotedblleft}, Flush, \"Odeme

Allgemeinsymptome: Juckreiz, Zittern, Schwindel, Kopfschmerz

\\\midrule
\gTabSub{II}

 &
RR ${\downarrow}$  , Tachykardie

\"Ubelkeit, Erbrechen, Durchfall 

\\\midrule
\gTabSub{III}
 &
Allgemeine Schocksymptome

Bronchospasmus
\\\midrule
\gTabSub{IV}
 &
Atem-/Kreislaufstillstand

Schwerer Schock
\\\bottomrule
\end{tabulary}
\end{center}

\begin{gMassnahmenEnv}
    \paragraph{Spezielle Ma{\ss}nahmen beim
    Anaphylaktischen Schock}
    
    \begin{itemize}
        \item Allgemein: \gEs{Allgemeine Ma{\ss}nahmen der
        Schockbek\"ampfung}
        (s.\,\ref{topic:schockbekaempfung})
        \item Antigen  nach M\"oglichkeit entfernen!
        \item Lagerung: situationsgerecht. Druckverlust im Vordergrund: Beine
        hoch.
        Atemprobleme im Vordergrund: stark erh\"ohter Oberk\"orper
    \end{itemize}
\end{gMassnahmenEnv}



\subsubsection{Septischer Schock und Toxischer Schock}
\label{topic:schock-septischer}Siehe auch: Kap.
\ref{topic:sepsis} im Hygiene-Teil 

\gDefinition{Schock infolge einer schweren, k\"orperweiten Infektion
oder aufgrund eines Gef\"a{\ss}-beeinflussenden Giftstoffes
(Toxins).}

Durch bakterielle Gef\"a{\ss}wandsch\"adigung, temperaturbedingt (hohes
Fieber {\textgreater}\,38{\textcelsius}) sowie aufgrund von
Giftstoffen kommt es zu Fl\"ussigkeitsaustritt in den Zwischen\-gewebs\-raum sowie
zu Gef\"a{\ss}weitstellung.

\paragraph{Ursachen}

Diverse Erreger und Toxine

\begin{itemize}
    \item Langanhaltender Entz\"undungsherd der pl\"otzlich Erreger
    {\quotedblbase}ins Blut streut{\textquotedblleft} und
    {\quotedblbase}explodiert{\textquotedblleft}
    \item Mit Keimen besiedelte Fremdk\"orper (diverse
    medizinische Sonden, Venflons, etc.)
    \item Tampons
    \item \ldots
\end{itemize}

\paragraph{Symptome}

\begin{itemize}
\item Schock
\item hohes Fieber oder kein Fieber
\item Haut hei{\ss}, trocken, rot oder bleich, grau,
{\quotedblbase}marmoriert{\textquotedblleft}
\end{itemize}

\begin{gMassnahmenEnv}
    \paragraph{Spezielle Ma{\ss}nahmen beim Septischen und
    Toxischen Schock}
    
    \begin{itemize}
        \item Allgemein: \gEs{Allgemeine Ma{\ss}nahmen der
        Schockbek\"ampfung} (s.\,\ref{topic:schockbekaempfung})
        \item Lagerung: situationsgerecht, Beine hoch
        \item Schutz vor Unterk\"uhlung/\"Uberw\"armung
        \item Aviso im Krankenhaus / Intensivstation
        ({\quotedblbase}\"Uberwachungsbett{\textquotedblleft})
        \item ben\"otigt Antibiotika (im KH)
    \end{itemize}
\end{gMassnahmenEnv}

\subsubsection{Neurogener Schock}
\gDefinition{Zusammenbruch der nerv\"osen Kreislaufregulation,
inad\"aquate arterielle und ven\"ose Gef\"a{\ss}erweiterung
\gZ{(Vasodilatation)}}

\paragraph{Ursachen} 

\begin{itemize}
    \item Neurologische Erkrankungen wie Querschnitt-Verletzung des
    R\"uckenmarks
    \item Die Gef\"a{\ss}e bekommen keine Nerven-Impulse und kontrahieren
    sich nicht mehr.
    \item ${\rightarrow}$ Blut versackt
\end{itemize}

\begin{gMassnahmenEnv}
    \paragraph{Spezielle Ma{\ss}nahmen}
    
    \begin{itemize}
        \item Allgemein: \gEs{Allgemeine Ma{\ss}nahmen der
        Schockbek\"ampfung} (s.\,\ref{topic:schockbekaempfung})
        \item Lagerung: Flachlagerung (WS-Verletzung ist Kontraindikation gg.
        Beine hoch!)
    \end{itemize}
\end{gMassnahmenEnv}
\section{Temperaturregulation}

W\"armeproduktion durch

\begin{itemize}
    \item Zellaktivit\"at
    \item Stoffwechselvorg\"ange
    \item Bewegung (Muskelatmung)
\end{itemize}
W\"armeabgabe \"uber:
\begin{itemize}
    \item Haut
    \item Schwei{\ss}
    
    \begin{itemize}
        \item Wirkung durch Verdunstungsk\"alte
        \item ca. 0,5\,l\,/\,Tag
        \item bis max. 1,5\,l\,/\,h bei schwerer
        k\"orperlicher  Arbeit
    \end{itemize}
\end{itemize}

Normale Körper-Kerntemperatur: 37\textcelsius


\section{Stoffwechsel}

\gDefinition{Abl\"aufe von Stoffumsetzung und 
Energiegewinnung, N\"ahrstoffe werden in W\"arme und Energie  umgewandelt}

\paragraph{Arten}

\begin{itemize}
    \item Zuckerstoffwechsel (Kohlenhydrate)
    \begin{itemize}
        \item Energielieferant
    \end{itemize}
    \item Fettstoffwechsel (Fetts\"auren, ges\"attigt und  unges\"attigt)
    \begin{itemize}
        \item Doppelt so hoher Energiegehalt wie  Kohlenhydrate
        \item L\"osungsmittel f\"ur Vitamine
        \item Bausteine der Zellmembran
        \item Bestandteil von Hormonen
    \end{itemize}
    \item Eiwei{\ss}stoffwechsel (Proteine)
    \begin{itemize}
        \item Zellmembranbausteine
        \item Enzyme!
    \end{itemize}
\end{itemize}

\section{Wasser- und Elektro\-lyt\-haus\-halt}
Das K\"orperwasser macht ca. 60\,\% des K\"orpergewichtes
aus und teilt sich auf:

\begin{itemize}
\item \gT{IntraZellularRaum} (IZR), 40\,\% KG 
\item \gT{ExtraZellularRaum} (EZR), 20\,\% KG 
\begin{itemize}
\item \index{Zwischengewebsraum}Zwischengewebsraum,
$\sim$\,16\,\%: {\quotedblbase}freier{\textquotedblleft}
Raum zwischen den Zellen im Gewebe
\item \index{Gef\"a{\ss}raum}Gef\"a{\ss}raum $\sim$\,4\,\%:
Raum in den Blutgef\"a{\ss}en
\end{itemize}
\end{itemize}

\paragraph{Trennung} 

\begin{itemize}
    \item IZR -- EZR:    Zellmembran (semipermeable Membran, siehe unten)
    \item Zwischengewebsraum -- Gef\"a{\ss}raum:  Blutgef\"a{\ss}wand
\end{itemize}


\paragraph{Atome, Ionen und Elektrolyte}
\index{Atome}\index{Ionen}\index{Elektrolyte}\label{topic:elektrolyte}

Atome sind die kleinsten Grundbestandteile von Marterie.

Ionen sind Atome die elektrisch geladen sind.

Geladene Ionen sind gut in Wasser l\"oslich. Wenn sie in Wasser gel\"ost
sind heissen sie
"`\gT{Elektrolyte}"'.

\gFazit{Im Wasser gel\"oste, elektrisch geladene Teilchen
hei{\ss}en \gT{Elektrolyte}.}

Die wichtigsten Ionen sind Natrium ($Na^+$), Kalium ($K^+$)
(positiv geladen) und Chlorid ($Cl^-$)(negativ geladen). 

Funktion:

\begin{itemize}
    \item \gEs{Verschiebung von
    Wasser} zwischen den einzelnen R\"aumen \gZ{(durch
    Osmose und Diffusion, siehe unten)}
    \begin{itemize}
        \item ${\rightarrow}$ Volumenregulation des K\"orpers, siehe unten.
    \end{itemize}
    \item \gEs{Elektrische Ladung}: Wie bei einer Batterie
    ist eine Zelle geladen (innen Minuspol, aussen Pluspol). Das ist besonders wichtig f\"ur
    \begin{itemize}
        \item Herz (Erregungsleitung!)
        \item Nerven
        \item Muskel
    \end{itemize}
\end{itemize}

\gZ{Elektrolyte sind geladene Teilchen, die  bei Aufspaltung
von S\"auren, Laugen  oder deren Salzen in w\"assriger L\"osung  entstehen.

\gE{Kationen:} positiv geladene Teilchen ($K^+$, $Na^+$,
$Mg^{++}$, $Ca^{++}$, \ldots); \gE{Anionen:} negativ
geladene Teilchen (Cl- )

Ionenverteilung: \gE{IZR:} viel $K^+$, wenig $Na^+$;
\gE{EZR:} viel $Na^+$, wenig $K^+$}



\subsection{Osmose und Diffusion} \index{Osmose}\index{Diffusion}

\gAuthNotes{GAB: Osmose und Diffusion: Grunds\"atzlich wichtig und
relevant, allenfalls bessere Erkl\"arung, aber bleibt in Normalschrift, nicht
heller. [GAB]

KOCH: Sollte alles grau werden}

\gMarried
{Geladene Teilchen k\"onnen die Zellmembran nicht selbstst\"andig
durchdringen, Wasser dagegen schon. Man nennt die Zellmembran daher
eine {\quotedblbase}\index{semipermeable
Membran}\gT{semipermeable Membran}{\textquotedblleft}. 

Die Natur m\"ochte dass in einer Fl\"ussigkeit \"uberall die gleiche
Konzentration von geladenen Teilchen herrscht. Da die Zellmembran (und
teilweise auch die Blutgef\"a{\ss}wand) nur das Wasser durchl\"asst,
nicht jedoch die geladenen Teilchen verschiebt sich nur das Wasser von
einem Raum in den anderen. Man kann sagen dass die Elekktrolyte einen
Druck auf das Wasser aus\"uben {\quotedblbase}die Seiten zu
wechseln{\textquotedblleft}. Man nennt diesen Druck den
{\quotedblbase}\index{osmotischer Druck}\gT{osmotischen
Druck}{\textquotedblleft}. Mann kann auch sagen
{\quotedblbase}das Wasser folgt dem Konzentrationsgef\"alle{\textquotedblleft}.}
{\begin{center}
\includegraphics[width=\linewidth]{./images/osmose.png}
\captionof{figure}{Diffusion. \gSourcePicture{Fischer}}
\end{center}}




\gZ{Daneben gibt es den {\quotedblbase}onkotischen
Druck{\textquotedblleft} der durch Eiwei{\ss}bestandteile ausge\"ubt wird.}

Wie kommen die Elektrolyte in die einzelnen R\"aume wenn doch nur Wasser
die Wand passieren kann?

Daf\"ur gibt es in der Zelle spezielle "`Pumpen"'
\gZ{(Transportproteine)} die die Ionen hin und her
bef\"ordern. Die Steuerung richtige Verteilung der verschiedenen Ionen ist recht kompliziert aber lebenswichtig. Eine {\quotedblbase}falsche{\textquotedblleft} Verteilung der Elektrolyte
kann schlimme Folgen haben (zB. Herzrhythmusst\"orungen). 

\gFazit{Die richtige Verteilung der Elektrolyte im K\"orper
ist lebenswichtig. Elektrolyst\"orungen k\"onnen lebensgef\"ahrlich sein. }

\paragraph{H\"aufige Ursachen von
Elektrolyst\"orungen}\index{Elektrolyst\"orungen}

\begin{itemize}
    \item Niereninsuffizienz (gest\"orte Ausscheidung)
    \begin{itemize}
        \item Besonders bei Dialysepatienten
    \end{itemize}
    \item Durchfall, Erbrechen (Verlust von Elektrolyten)
    \item Medikamente
    \item Vergiftung
\end{itemize}



\subsection{Wasser- und Elektrolythaushalt}

Blutvolumen = Plasmawasser + zellul\"are Bestandteile (Blutk\"orperchen)

Beim Erwachsenen 7-8\,\% des Körpergewichts

Beim Neugeborenen / S\"augling 8-9\,\% des Körpergewichts  

Verluste bis 20\,\% ohne Funktionsbeeintr\"achtigung
tolerierbar, allerdings Sch\"aden in der  Mikrozirkulation

${\rightarrow}$ auch ohne klinische Anzeichen ist bei
Volumensverlust ein Volumensersatz erforderlich!

\begin{itemize}
    \item osmotischer Druck
    \begin{itemize}
        \item Volumenregulation des K\"orpers
    \end{itemize}
    \item Konzentrationsgef\"alle ${\rightarrow}$ Spannungsunterschied an
    den Zellmembranen ${\rightarrow}$ Grundlage f\"ur Erregungsleitung /
    Muskelkontraktion
    \item Regelung durch:
    \begin{itemize}
        \item Hormone \gZ{(Renin/Angiotensin/Aldosteron)}
        \item Ausscheidung \"uber Niere
    \end{itemize}
\end{itemize}

\subsection{S\"aure-/Basenhaushalt}


\index{S\"aure}\gT{S\"auren} = Verbindungen, die 
Wasserstoffionen ($H^+$) abgeben k\"onnen

\index{Lauge}\gT{Laugen} = \index{Base}\gT{Basen} =
Verbindungen, die $H^+$ aufnehmen k\"onnen

\paragraph{\gT{pH}-Wert}\index{pH-Wert}

\gDefinition{Der pH-Wert gibt an wie sauer oder basisch eine
L\"osung ist.}

pH\,{\textless}\,7 ist sauer, pH\,=\,7 ist neutral,
pH\,{\textgreater}\,7 ist basisch

\gZ{pH normal ca. um 7,4.

pH {\textless} 7,35  $\rightarrow$ Azidose, pH {\textgreater} 7,45  $\rightarrow$ Alkalose}

Der richtige pH-Wert ist lebenswichtig. Der K\"orper ist da sehr
empfindlich, der pH-Wert muss sich in engen grenzen bewegen. Schon
geringe Abweichungen k\"onnen gro{\ss}e Folgen haben. 

Leider gibt es im Rettungsdienst normalerweise keine M\"oglichkeit den
pH-Wert zu messen, jedoch haben wir gro{\ss}en Einflu{\ss} auf ihn
durch unsere Ma{\ss}nahmen, siehe unten. 

Bei einer \"Ubers\"auerung sprechen wir von einer
\index{Azidose}\gT{Azidose}, wenn der Patient basisch wird
von einer \index{Alkalose}\gT{Alkalose}.

\gFazit{Bei einer \gEs{\"Ubers\"auerung} sprechen wir von
einer \gT{Azidose}, wenn der Patient \gEs{basisch} wird
von einer \gT{Alkalose}.}

\paragraph{Regulationsmechanismen}

\begin{itemize}
    \item Blut mit Puffersystem
    \begin{itemize}
        \item \gZ{$H^+ + HCO_3^- \leftrightarrow H_2O + CO_2$  (Formel nicht lernen!)}
        \item S\"aure {\tiny + Bikarbonat}
        ${\leftrightarrow}$ {\tiny Wasser +} $CO_2$
        \begin{itemize}
            \item Atmet der Patient nicht (genug) staut sich das CO2 und bildet mit
            Wasser $H^+"$ ${\rightarrow}$ Atmungsbedingte
            \"Ubers\"auerung
            \item Atmet der Patient zu viel wird zu viel CO2 abgeatmet und wird von
            H+ und dem anderen Stoff nachgebildet ${\rightarrow}$ H+ fehlt
            ${\rightarrow}$ Atmungsbedingte Alkalose
            \item F\"allt zu viel $H^+$ (S\"aure) im
            K\"orper an (durch den Stoffwechsel, Stoffwechselbedingte \"Ubers\"auerung) wird es zu (Wasser
            und) $CO_2$ umgewandelt, das $CO_2$ wird
            abgeatmet
        \end{itemize}
    \end{itemize}
    \item Lunge
    \begin{itemize}
        \item Abatmung von $CO_2$ ${\rightarrow}$ Einfluss auf den Puffer (s.o.)
    \end{itemize}
    \item Niere
    \begin{itemize}
        \item Ausscheidung (oder Wiederaufnahme)  von H+
    \end{itemize}
\end{itemize}

\subsubsection{Atmungsbedingte \"Ubers\"auerung - Respiratorische Azidose}
\label{bkm:RefHeading35233321}\index{\"Ubers\"auerung}\index{Respiratorische
Azidose}durch Hypoventilation / verminderten  Gasaustausch:

$CO_2$  Abatmung vermindert ${\rightarrow}$ Anhäufung von
$CO_2$ im Blut\footnote{\gZ{Hyperkapnie}}, Ans\"auerung des
Blutes

\paragraph{Ursachen}

\begin{itemize}
    \item Insuffiziente Atmung / Atemstillstand
\end{itemize}

\paragraph{Therapie}

\begin{itemize}
    \item Erh\"ohung des Atem-Minuten-Volumens
\end{itemize}


\subsubsection{Stoffwechselbedingte \"Ubers\"auerung
-Metabolische Azidose} \index{\"Ubes\"auerung}\index{Metabolische Azidose}Zunahme von S\"auren im  K\"orper

\paragraph{Ursachen}

\begin{itemize}
    \item Verminderte Nierenleistung
    \item Schock
    \item Durchfall
    \item Diabetisches Koma
\end{itemize}

\paragraph{Gefahren}

\begin{itemize}
    \item Sch\"adigung des Zellstoffwechsels
    \item Sch\"adigung des Herzmuskels
\end{itemize}

\subsubsection{Atmungsbedingte (respiratorische) Alkalose}
\index{Alkalose}Hyperventilationssyndrom

Elektrolytverschiebung: Kribbeln in den Fingern, Pf\"otchenstellung,
Kr\"ampfe

\subsubsection{Stoffwechselbedingte (metabolische) Alkalose}
z.B. durch anhaltendes Erbrechen: Verlust von S\"aure

\subsubsection*{Zusammenfassung}

\gFazit{Wer zu wenig atmet wird sauer.

Wer zu viel atmet wird basisch und sp\"urt ein Kribbeln in den Fingern. 

Wer Probleme mit dem Stoffwechsel hat kann auch sauer oder basisch
werden.}








\section{Tod}

Der Tod ist eine der grundlegensten und unausweichlischsten
Vitalfunktionen überhaupt.  



\paragraph{Klinischer Tod vs. biologischer Tod}

\begin{center}
\begin{tabularx}{\linewidth}{XX}
\toprule
\gTabHeading{Klinischer Tod}
 &
\gTabHeading{Biologischer Tod}
\\\midrule
keine Atmung

kein Puls

Reanimation mu{\ss} innerhalb der ersten  Minuten erfolgen

(dann unter Umst\"anden) reversibel
 &
Organtod des Gehirns (irreversibel)

gleichzusetzen mit Tod des Individuums
\\\bottomrule
\end{tabularx}
\end{center}

\paragraph{Wann ist es wirklich zu sp\"at?}

\gFazit{Todesfeststellung: Grunds\"atzlich NUR durch den
Arzt!}

\gAuthNotes{Todeszeichen: Nicht wie von Roman gefordert in
Tabelle mit mit sicheren und unsicheren Todeszeichen umgebaut, halte ich f\"ur un\"ubersichtlich
und bef\"urchte Verwechslungsgefahr mit
Knochenbruchzeichen.[GAB]}





\subsubsection{Zeichen des Todes und Leichenzeichen}

\gZ{nach: \cite{ForMed2}}

Grundsätzlich gibt es 3 Zeichen die als \gEs{sichere
Todeszeichen} gelten:

\begin{enumerate}
    \item \gEs{Bildung von blauvioletten Totenflecken} durch
    Absackung des Blutes an die untenliegenden Stellen.
    \item \gEs{Nicht mit dem Leben vereinbare Verletzungen}
    wie zum Beispiel eine Abtrennung des Kopfes vom Köper.
    \item \gEs{Fäulniserscheinungen}
\end{enumerate}

Die ersten beiden Punkte sind \gE{frühe
Leichenerscheinungen}, da sie relativ rasch auftreten.
Andere frühe Leichenerscheinungen sind die \gE{Abkühlung
der Leiche} und die \gE{Totenstarre}. 

Längere Zeit nach dem Tod treten die \gE{späten
Leichenerscheinungen} auf, dazu gehören \gE{Fäulnis und
Verwesung}, \gE{Mumifizierung}, \gE{Tierfraß} und
\gE{Skelettierung}. 

Die Feststellung der Todeszeichen und Leichenerscheinungen
ist jedoch oft nicht so einfach, die Anzeichen können oft
auch mit anderen Sachen verwechselt werden, z.B.
Totenflecke mit Blutergüssen oder die Totenstarre mit einem
tonischen Krampf.

\gFazit{Todesfeststellung: Grunds\"atzlich NUR durch den
Arzt!}

\gCave{Nobody is dead until warm and dead!

Niemand wird für tot erklärt, bevor er nicht erwärmt
wurde.}

Unter den folgenden Voraussetzungen darf eine
Reanimation unterlassen werden:

\begin{itemize}
    \item \gEs{Verwesung}
    \item \gEs{Fäulnis}
    \item \gEs{Mumifizierung}
    \item \gEs{Skelettierung}
    \item \gEs{Verletzungen, die nicht mit dem Leben 
    vereinbar sind}
    \item \gEs{Tierfraß}
    \item Verbindliche oder Beachtliche
    \gEs{Patientenverfügung} wie unter
    \ref{topic:patientenverfuegung} ausgefhrt
    \item Abbruch der Reanimationsma{\ss}nahmen aufgrund
    von \gEs{Ersch\"opfung}
\end{itemize}

\putbib




